\section{Introduction (Введение)}
\label{sec:introduction}

\subsection{Индустриальный контекст и дилемма скипов}
Современные рекомендательные системы музыкальных стриминговых сервисов (Яндекс Музыка, Spotify, Apple Music) оперируют в условиях принципиально непрерывного и неструктурированного потребления контента. В отличие от платформ электронной коммерции, где транзакция (покупка) однозначно подтверждает позитивный интерес, в потоковом аудио взаимодействие пользователя носит бимодальный характер~\cite{joachims2005accurately,jannach2021recommender}.

Анализ логов взаимодействия на масштабном датасете Yandex Music (\textbf{Yambda-50M})~\cite{yandex2024yambda} показывает, что от \textbf{30\% до 35\%} всех зафиксированных событий представляют собой мгновенные скипы (время непрерывного воспроизведения $D < 30$ секунд или $0.0$ секунд, когда пользователь переключает трек в первые секунды из-за дискомфорта или рассогласования настроения).

\begin{figure}[htbp]
    \centering
    \includegraphics[width=0.88\textwidth]{fig1_the_skip_dilemma.png}
    \caption{\textbf{Дилемма скипов в логах потокового стриминга.} Распределение длительности прослушивания носит выраженный бимодальный характер: 35\% событий приходятся на короткие скипы ($<30$ с), тогда как 65\% отражают вовлечённое прослушивание композиции. Оптимизация правдоподобия всех наблюдаемых переходов без разделения качества наказывает модель за отказ от скипов.}
    \label{fig:skip_dilemma}
\end{figure}

\subsection{Проблема False Negatives при оптимизации правдоподобия (MLE)}
Классическая парадигма обучения последовательных рекомендательных систем (SASRec~\cite{kang2018self}, BERT4Rec~\cite{sun2019bert4rec}, Causal Transformers~\cite{radford2019language}) сводится к задаче предсказания следующего токена (\textit{Next-Item Prediction}) через минимизацию кросс-энтропии:
\begin{equation}
    \mathcal{L}_{\text{MLE}}(\theta) = - \sum_{t=1}^{T} \log \pi_\theta(a_t \mid s_{<t}),
    \label{eq:mle_loss}
\end{equation}
где $s_{<t} = (a_1, \dots, a_{t-1})$ — предыстория пользователя, а $a_t$ — фактически зафиксированный в логах следующий трек.

В этой постановке возникает фундаментальное концептуальное противоречие:
\begin{enumerate}
    \item \textbf{Штраф за правильный отказ (False Negative):} Если пользователь включил трек $a_t$ и переключил его через 1 секунду из-за раздражения, функция потерь $\mathcal{L}_{\text{MLE}}$ всё равно штрафует модель, если та присвоила треку $a_t$ низкую апостериорную вероятность.
    \item \textbf{Слепота к удовлетворённости:} Модель учится воспроизводить историю кликов, а не максимизировать качество потребления (время удержания, дослушивания до конца).
    \item \textbf{Эксплуатация безопасных поп-хитов:} Модель сваливается в гипер-рекомендацию популярных треков с высокой априорной частотой встречаемости, снижая разнообразие каталога.
\end{enumerate}

\subsection{Исследовательские вопросы (Research Questions)}
Для разрешения дилеммы скипов в данной работе исследуется адаптация методов обучения с подкреплением и согласования предпочтений (RL Alignment) для причинного языкового трансформера. Мы формулируем пять ключевых исследовательских вопросов:
\begin{itemize}
    \item \textbf{RQ1 (Поведение политики и динамика):} Каким образом оптимизация маржи предпочтений (DPO) перестраивает вероятностное распределение модели по сравнению с базовой MLE-политикой, и как быстро происходит сходимость скрытого вознаграждения?
    \item \textbf{RQ2 (Налог на выравнивание и компромисс Парето):} Приводит ли оптимизация вознаграждения к деградации ранжирования следующего трека (Alignment Tax), и способна ли вспомогательная регуляризация $\mathcal{L}_{\text{SFT}}$ устранить этот компромисс?
    \item \textbf{RQ3 (Разнообразие каталога и монополия хитов):} Схлопывается ли RL-политика в узкий набор безопасных треков, или, напротив, стимулирует обнаружение качественного контента из длинного хвоста (\textit{Long-Tail Discovery})?
    \item \textbf{RQ4 (Авторегрессионные роллауты и Exposure Bias):} Как ведёт себя выровненная модель при формировании очереди воспроизведения из 5 последовательных треков без обратной связи в условиях накопления ошибки авторегрессии?
    \item \textbf{RQ5 (Долгосрочные прокси и онлайн A/B валидация):} Как связать сессионные суррогаты вознаграждения с долгосрочным удержанием (Retention D7/D30) в дизайне промышленного эксперимента?
\end{itemize}

\subsection{Научные и инженерные вклады работы}
В рамках проведённого исследования получены следующие ключевые результаты:
\begin{enumerate}
    \item \textbf{Мультипарадигмальный сравнительный анализ:} На едином индустриальном датасете Yambda-50M сопоставлены четыре парадигмы: базовый Causal GPT-2, 2-стадийное реранжирование с функцией ценности $Q(s,a)$, Actor-Critic PPO~\cite{schulman2017proximal} и сквозной Direct Preference Optimization (DPO)~\cite{rafailov2023direct}.
    \item \textbf{Концепция Zero Alignment Tax:} Доказано, что добавление SFT-регуляризации ($\gamma=0.10$) предотвращает распад семантической структуры рекомендаций, сохраняя и увеличивая точность ранжирования (+0.67\% HR@10, $p < 0.0001$).
    \item \textbf{Эмпирический аудит 3 стратегий вознаграждения:} Исследованы бинарный порог ($<30$ с), непрерывная доля дослушивания ($D/T$) и композитная награда с защитой от повторов, выявившие преимущество порогового отсечения шума.
    \item \textbf{Оздоровление каталога:} Показано, что DPO снижает корреляцию с популярностью на 12.1\% и открывает пользователям более 1\,000 новых треков из длинного хвоста.
\end{enumerate}
